\documentclass[11pt,a4paper]{article}

\usepackage[a4paper,margin=1in]{geometry}
\usepackage{booktabs}
\usepackage{longtable}
\usepackage{hyperref}
\usepackage{enumitem}
\usepackage{graphicx}
\usepackage{xcolor}
\usepackage{array}

\title{
EcoTwin \
Product Requirements Document (PRD)
}

\author{
Ecosystem Analog Discovery and Forecasting Platform
}

\date{\today}

\begin{document}

\maketitle
\tableofcontents
\newpage

% ======================================================
\section{Document Information}

\begin{tabular}{|p{4cm}|p{8cm}|}
\hline
Project Name & EcoTwin \
\hline
Version & 1.0 \
\hline
Document Type & Product Requirements Document \
\hline
Prepared By & Project Team \
\hline
Target Release & MVP Release \
\hline
\end{tabular}

\vspace{1cm}

% ======================================================
\section{Executive Summary}

EcoTwin is an ecosystem discovery and forecasting platform that enables users to explore ecosystems using satellite imagery, identify ecologically similar regions, analyze historical environmental trends, and generate future ecosystem forecasts.

The platform leverages Earth Observation data, geospatial foundation models, ecosystem similarity search, and analog forecasting to help researchers, environmental planners, policymakers, and conservation organizations better understand ecosystem evolution.

The system will provide an intuitive web-based interface where users can interactively select geographic locations, discover ecosystem analogs, visualize environmental changes, and generate automated ecosystem reports.

% ======================================================
\section{Problem Statement}

Environmental ecosystems are continuously changing due to climate variability, urbanization, agricultural expansion, and natural processes.

Current monitoring systems typically focus on descriptive observations and rarely provide mechanisms to:

\begin{itemize}
\item Discover ecologically similar ecosystems.
\item Compare ecosystem trajectories.
\item Learn from historical analog ecosystems.
\item Forecast future ecosystem changes.
\end{itemize}

Researchers and planners often lack tools that can answer:

\begin{quote}
``What ecosystems elsewhere have evolved similarly to this ecosystem, and what can we learn from their historical evolution?''
\end{quote}

EcoTwin aims to address this challenge.

% ======================================================
\section{Vision}

To build an intelligent ecosystem intelligence platform capable of:

\begin{itemize}
\item Discovering ecosystem analogs.
\item Understanding ecosystem evolution.
\item Providing explainable ecosystem forecasts.
\item Supporting environmental decision making.
\end{itemize}

% ======================================================
\section{Goals}

\subsection{Business Goals}

\begin{enumerate}
\item Demonstrate ecosystem analog search.
\item Enable interactive ecosystem exploration.
\item Support environmental forecasting workflows.
\item Provide explainable ecosystem insights.
\end{enumerate}

\subsection{Product Goals}

\begin{enumerate}
\item Allow users to select any location within the study area.
\item Automatically identify the corresponding ecosystem region.
\item Retrieve similar ecosystem regions.
\item Display ecosystem evolution through time.
\item Generate analog-based forecasts.
\item Generate downloadable ecosystem reports.
\end{enumerate}

% ======================================================
\section{Target Users}

\subsection{Primary Users}

\begin{itemize}
\item Environmental Researchers
\item Conservation Scientists
\item GIS Analysts
\item Remote Sensing Researchers
\item Ecologists
\end{itemize}

\subsection{Secondary Users}

\begin{itemize}
\item Policymakers
\item NGOs
\item Government Agencies
\item Academic Institutions
\item Students
\end{itemize}

% ======================================================
\section{User Personas}

\subsection{Persona 1: Environmental Researcher}

\textbf{Name:} Research Scientist

Needs:

\begin{itemize}
\item Discover similar ecosystems.
\item Analyze ecosystem evolution.
\item Generate scientific insights.
\end{itemize}

Success Criteria:

\begin{itemize}
\item Can retrieve ecosystem analogs.
\item Can compare environmental trajectories.
\end{itemize}

\subsection{Persona 2: Conservation Planner}

Needs:

\begin{itemize}
\item Understand ecosystem risks.
\item Forecast future ecosystem conditions.
\item Support planning decisions.
\end{itemize}

Success Criteria:

\begin{itemize}
\item Receives explainable forecasts.
\item Can generate ecosystem reports.
\end{itemize}

% ======================================================
\section{Product Scope}

\subsection{In Scope}

\begin{itemize}
\item Interactive ecosystem map.
\item Location selection.
\item Ecosystem region identification.
\item Ecosystem similarity search.
\item Historical trend visualization.
\item Analog forecasting.
\item Automated reporting.
\end{itemize}

\subsection{Out of Scope}

\begin{itemize}
\item Real-time satellite ingestion.
\item Species-level prediction.
\item Climate simulation.
\item Hydrological simulation.
\item Global-scale deployment.
\end{itemize}

% ======================================================
\section{User Journey}

\begin{enumerate}
\item User opens EcoTwin.
\item User navigates the map.
\item User clicks a location.
\item System identifies the ecosystem region.
\item System displays ecosystem summary.
\item User requests similar ecosystems.
\item System retrieves ecosystem analogs.
\item User explores historical trajectories.
\item User generates ecosystem forecast.
\item User downloads report.
\end{enumerate}

% ======================================================
\section{User Stories}

\subsection{Location Selection}

\textbf{US-001}

As a user,

I want to select a location on the map,

So that I can analyze the ecosystem around that location.

\vspace{0.5cm}

Acceptance Criteria:

\begin{itemize}
\item User can click on a map.
\item Selected region is highlighted.
\item Region metadata is displayed.
\end{itemize}

% ------------------------------------------------------

\subsection{Ecosystem Identification}

\textbf{US-002}

As a user,

I want the system to identify ecosystem characteristics,

So that I can understand the selected region.

Acceptance Criteria:

\begin{itemize}
\item Ecosystem type displayed.
\item Environmental indicators displayed.
\item Region summary generated.
\end{itemize}

% ------------------------------------------------------

\subsection{Similarity Search}

\textbf{US-003}

As a user,

I want to discover similar ecosystems,

So that I can compare ecological behavior.

Acceptance Criteria:

\begin{itemize}
\item Top similar ecosystems returned.
\item Similarity score displayed.
\item Results shown on map.
\end{itemize}

% ------------------------------------------------------

\subsection{Historical Analysis}

\textbf{US-004}

As a user,

I want to view historical ecosystem evolution,

So that I can understand long-term changes.

Acceptance Criteria:

\begin{itemize}
\item Historical charts available.
\item Multiple years visualized.
\item Comparison between ecosystems supported.
\end{itemize}

% ------------------------------------------------------

\subsection{Forecasting}

\textbf{US-005}

As a user,

I want ecosystem forecasts,

So that I can anticipate future changes.

Acceptance Criteria:

\begin{itemize}
\item Forecast generated.
\item Confidence score displayed.
\item Forecast explanation provided.
\end{itemize}

% ------------------------------------------------------

\subsection{Report Generation}

\textbf{US-006}

As a user,

I want downloadable reports,

So that I can share ecosystem insights.

Acceptance Criteria:

\begin{itemize}
\item PDF generated.
\item Charts included.
\item Analog ecosystems included.
\item Forecast included.
\end{itemize}

% ======================================================
\section{Functional Requirements}

\subsection{FR-001}

The system shall allow users to select locations on an interactive map.

\subsection{FR-002}

The system shall identify the ecosystem region associated with the selected location.

\subsection{FR-003}

The system shall display ecosystem metadata.

\subsection{FR-004}

The system shall retrieve similar ecosystems.

\subsection{FR-005}

The system shall display similarity scores.

\subsection{FR-006}

The system shall display ecosystem analogs on a map.

\subsection{FR-007}

The system shall visualize historical environmental indicators.

\subsection{FR-008}

The system shall support comparison of multiple ecosystem trajectories.

\subsection{FR-009}

The system shall generate ecosystem forecasts.

\subsection{FR-010}

The system shall generate downloadable reports.

% ======================================================
\section{Non-Functional Requirements}

\subsection{Performance}

\begin{itemize}
\item Region lookup < 1 second
\item Similarity search < 2 seconds
\item Forecast generation < 5 seconds
\item Report generation < 10 seconds
\end{itemize}

\subsection{Scalability}

System should support:

\begin{itemize}
\item 5,000+ ecosystem regions
\item Historical records for multiple years
\item Concurrent user access
\end{itemize}

\subsection{Usability}

\begin{itemize}
\item Intuitive interface
\item Minimal learning curve
\item Interactive visualizations
\end{itemize}

% ======================================================
\section{Success Metrics}

\subsection{Product Metrics}

\begin{itemize}
\item Successful location selection rate > 95%
\item Similarity retrieval success > 90%
\item Report generation success > 95%
\end{itemize}

\subsection{Performance Metrics}

\begin{itemize}
\item Average search latency < 2 seconds
\item Average forecast generation < 5 seconds
\item Average page load < 3 seconds
\end{itemize}

\subsection{User Experience Metrics}

\begin{itemize}
\item User task completion rate > 90%
\item Positive usability feedback
\item Successful ecosystem discovery workflows
\end{itemize}

% ======================================================
\section{Risks}

\begin{tabular}{|p{5cm}|p{4cm}|p{4cm}|}
\hline
Risk & Impact & Mitigation \
\hline
Satellite Data Quality &
Medium &
Cloud filtering \
\hline
Forecast Uncertainty &
High &
Confidence Scores \
\hline
Large Data Volume &
Medium &
Caching and indexing \
\hline
Similarity Retrieval Quality &
High &
Embedding validation \
\hline
\end{tabular}

% ======================================================
\section{Release Plan}

\subsection{Phase 1}

Core Infrastructure

\begin{itemize}
\item Map Interface
\item Region Selection
\item Ecosystem Summary
\end{itemize}

\subsection{Phase 2}

Ecosystem Discovery

\begin{itemize}
\item Similarity Search
\item Analog Retrieval
\item Historical Visualization
\end{itemize}

\subsection{Phase 3}

Forecasting and Reporting

\begin{itemize}
\item Analog Forecasting
\item PDF Reports
\item Final Optimization
\end{itemize}

% ======================================================
\section{Conclusion}

EcoTwin will provide a comprehensive ecosystem discovery and forecasting platform that combines ecosystem analog search, temporal analysis, and explainable forecasting into a unified user experience.

The platform will enable researchers, planners, and environmental stakeholders to better understand ecosystem evolution and support evidence-based environmental decision making.

\end{document}
